\documentclass[10pt,a4paper]{amsart}
\usepackage[margin=19mm]{geometry}
\usepackage{amsmath,amssymb,mathtools}
\usepackage[scr=rsfs,cal=euler]{mathalfa}
\usepackage{xfrac}
\usepackage[unicode,hidelinks,bookmarks=false]{hyperref}
\newcommand{\ie}{i.e.\ }
\newcommand{\cf}{cf.\ }
\newcommand{\resp}{respectively\ }
\newcommand{\Subsection}{\subsection}
\providecommand{\nobreakdash}{\nobreak\hbox{-}}
\setlength{\emergencystretch}{2em}
\multlinegap0pt
\usepackage{bm}
\usepackage{bbm}
\usepackage{dsfont}
\usepackage{xspace}
\usepackage{cancel}
\usepackage{mathtools}
\usepackage{amsthm}
\makeatletter
\@ifundefined{Proposition}{\newtheorem{Proposition}{Proposition}}{}
\makeatother
\title[Euclidean distance geometry and the orthogonal beltway probl]{Euclidean distance geometry and the orthogonal beltway problem}
\author{Dan Edidin, Arun Suresh}
\date{Source: arXiv:2604.26030v1 (CC BY 4.0)}
\begin{document}
\maketitle

\noindent\emph{Source and licence.} Euclidean distance geometry and the orthogonal beltway problem. Version arXiv:2604.26030v1 (CC BY 4.0), distributed under the Creative Commons Attribution 4.0 International licence, as recorded in the arXiv licence metadata for this exact version and shown on the abstract page https://arxiv.org/abs/2604.26030v1. Full rights notice: licenses/arxiv-2604.26030-CC-BY-4.0.txt.\par\smallskip

\noindent\textit{Editorial scope.} Counted unit: the Proposition whose source label is \texttt{prop:cayleymengerprop} in the article (source0.tex lines 1374--1377, source label \texttt{prop:cayleymengerprop}), delivered together with the article's own complete original proof of it (source0.tex lines 1378--1411, one \texttt{proof} environment of 34 lines). No mathematics is written, repaired or re-derived: statement and proof are the article's own text line for line. Required context retained uncounted: none. Every definition, display and label of those lines is retained unchanged. Omitted from this package and not redistributed: 1--1373, 1412--1412. Nothing inside a retained excerpt is omitted or altered: every retained body is delivered exactly as the article's lines are written, so no omission receipt of any kind accompanies this package. Citations: the retained excerpts carry no citation command at all, so no bibliography entry of the article is transcribed and no reference is redistributed. Reference adaptation: none was needed and none was made; every reference inside the delivered unit resolves inside it, so no unresolved reference or citation is produced. Printed slips kept literal: no printed word, symbol, hypothesis or number of the counted statement or of its proof was altered. Layout and class: the article's own class is \texttt{article} with packages including inputenc, footmisc, authblk, geometry, amsmath, amssymb, amsthm, stmaryrd, graphicx, subcaption, algorithm2e, caption, tikz, bbm; the wrapper is the lane's pinned \texttt{amsart} editorial wrapper at 10pt on A4, which restates the theorem-like environments the retained text needs and adds the article's own macro definitions that the retained text needs (none), with extra wrapper packages (bm, bbm, dsfont, xspace, cancel, mathtools). The delivered numbering is the unit's own. Assets: no retained span contains a figure environment, a graphics-inclusion command, a PDF-page inclusion, a table or a TikZ picture; no image file is copied into this package, the package declares no asset dependency and no image file is redistributed with it. Licence: version arXiv:2604.26030v1 of this article is distributed under the Creative Commons Attribution 4.0 International licence, as recorded in the arXiv licence metadata for this exact version and shown on the abstract page https://arxiv.org/abs/2604.26030v1. A full rights notice is delivered at licenses/arxiv-2604.26030-CC-BY-4.0.txt. Characters: every retained excerpt is pure ASCII, so the delivered engine, XeLaTeX with its default Latin Modern text font, reports no missing character.
\begin{Proposition}\label{prop:cayleymengerprop}
   Let $x,y,r$ be fixed distinct positive numbers such that $r \neq 1$ and $|r-1| < x,y < r+1$.
   Then the biquadratic polynomial $f(z) = D(1,1,r,x,y,z)$  has four distinct real roots.
\end{Proposition}

\begin{proof}
   Given $E = (1,1,r,x,y,z)$ with $r,x,y$  fixed then $D(E) = \alpha z^4 + \beta z^2 + \gamma$
    where
   \begin{equation*}
       \begin{split}
           \alpha(x,y) &= -2r^2\\
           \beta(x,y) &= -2+4r^2-2r^4+2x^2+2r^2x^2+2y^2+2r^2y^2-2x^2y^2\\
           \gamma(x,y) &= -2x^4 + 4x^2y^2 - 2y^4
       \end{split}
   \end{equation*}
   Substituting $t=z^2$ into $D(E)$ it suffices to show that the polynomial  $D_E(t) = \alpha t^2+\beta t+\gamma c$  has two positive real roots $t_1$ and $t_2$. \\

    We begin by analyzing the signs of the coefficients $\alpha,\beta$ and $\gamma$. Certainly $\alpha < 0$ and $\gamma = -2(x^2-y^2)^2 < 0$. 
    %To determine the sign of $\beta$ we first observe that for $i=1,2$, $$|r-1| = |\|v\|-\|v_i\|| \leq \|v_i-v\| \leq \|v_i\|+\|v\| = 1+r$$ which forces $(x,y) \in D(r) = [|r-1|, r+1]^2$ for any fixed $r>0$. 
    We will show that when $r \neq 1$, $\beta (x,y) \geq 0$ when $|r-1| < x,y < r+1$, by minimizing $\beta$ as a function of $x, y$ over the
    the square $[|r-1|,r+1] \times [|r-1|, r+1]$ and showing that the minimum is non-negative. To this end note that 
    $\nabla \beta = \langle 4x(1+r^2-y^2), 4y(1+r^2-x^2)\rangle$ vanishes if and only if $x=y=0$ or $x=y=\sqrt{1+r^2}$. However, since $r \neq 1$ and $x,y > |r-1|$, we obtain $x,y\neq 0$, forcing $x=y=\sqrt{1+r^2}$ to be the only relevant critical point in the interior of 
    the square $[|r-1|, r+1] \times [|r-1|, r+1]$. Any other extrema of $\beta(x,y)$ over the square  occur at its boundary. A similar verification along each edge of the square shows that there are no non-trivial critical points along the boundary, leaving the extreme points along the boundary to occur at its vertices. Evaluating $\beta(x,y)$ at its unique internal critical point and at each of the four vertices of the square
    \begin{equation*}
        \begin{split}
            \beta(\sqrt{1+r^2}, \sqrt{1+r^2}) = r^2(1+r^2) &> 0\\
            \beta(r+1,r+1) = \beta(|r-1|,|r-1|) &= 0\\
            \beta(|r-1|,r+1) = \beta(r+1,|r-1|) = 16r^2 &> 0
        \end{split}
    \end{equation*}
    Thus, by the extreme value theorem, $\beta(x,y) >0$ on the square when $|r-1| < x,y < r+1$ as desired. 
    

Having determined the signs of $\alpha,\beta$ and $\gamma$ - we can derive the desired result via the following observation: the roots $t_1$ and $t_2$, if they exist, must satisfy the Vieta relations $$t_1t_2 = \gamma/\alpha > 0 \text{ and } t_1 + t_2 = -\frac{\beta}{\alpha} \geq 0.$$ From the first equation, we deduce that if $t_1, t_2$ are real then they have the same sign. From the second equation, we deduce that this sign has to be positive. Finally, a routine computation ensures that the discriminant $\Delta$ of $D_E(t)$ is 
positive on the interior of the square $[|r-1|, r+1] \times [|r-1|, r+1]$ so
the two roots $t_1$ and $t_2$ are indeed real and positive Hence the biquadratic polynomial $f(z)$ has two positive real roots.    
   
%    Moreover, the discriminant $\Delta$ only vanishes at the extreme values $(x,y) = (|r-1|, |r-1|)$ and $(x,y) = (r+1,r+1)$ which correspond to $v_1 = v_2$. Thus $\Delta > 0$ in all regions of interest, and as a result, $D(E)$ has four distinct real roots. \arun{[[Arun: Though this doesn't happen in practice, in the extreme case where the discriminant is zero, we have $t_1 = t_2$ and therefore $z$ (corresponding to the missing edge length) is uniquely determined. As such, we can remove this line and instead say in the statement of the proposition ``the biquadratic polynomial $f(z) = D(1,1,r,x,y,z)$ has two positive real roots (counted with multiplicity)".]]}
\end{proof}

\end{document}
